\documentclass[aspectratio=169,8pt]{beamer}
\usetheme{Madrid}
\usepackage[utf8]{inputenc}
\usepackage[T1]{fontenc}
\usepackage{amsmath,amssymb}
\usepackage{booktabs}
\usepackage{enumitem}
\setlist[enumerate]{label=\Alph*.,leftmargin=*,itemsep=1pt,topsep=2pt}
\setlist[itemize]{leftmargin=*,itemsep=1pt,topsep=2pt}
\setbeamerfont{block body}{size=\tiny}
\setbeamerfont{block title}{size=\scriptsize}
\setbeamerfont{frametitle}{size=\footnotesize}
\setbeamerfont{normal text}{size=\tiny}
\setbeamertemplate{navigation symbols}{}
\setbeamertemplate{itemize item}{\tiny$\bullet$}
\setbeamertemplate{itemize subitem}{\tiny$\circ$}

\title[GARP RAI]{GARP Risk \& AI --- Q551--Q650}
\subtitle{Unofficial Exam Prep}
\author{Unofficial}
\date{\today}

\begin{document}
	
	\begin{frame}
		\titlepage
	\end{frame}
	
	% =================================================================
	\section{Advanced AI Risk Taxonomy}
	% =================================================================
	
	\begin{frame}{Q557 --- Model Inventory vs AI Portfolio}
		\begin{block}{Question}
			How does an AI portfolio differ from a model inventory?
			\begin{enumerate}
				\item The portfolio and inventory are functionally isomorphic artefacts, both maintaining versioned metadata schemas, ownership attestations, and validation-state enumerations across the model estate.
				\item The portfolio is a strategic capability map spanning build-versus-buy decisions, capital allocation, and value-realisation roadmaps, whereas the inventory is an operational registry of deployed artefacts with lineage, drift telemetry, and lifecycle-state metadata.
				\item The portfolio is scoped exclusively to third-party procured solutions governed by vendor risk schedules, while the inventory is restricted to internally developed models that have entered post-retirement archival states.
				\item The inventory is a forward-looking capability roadmap with investment theses, while the portfolio is a real-time inference-call ledger capturing latency percentiles, throughput counters, and token-consumption metrics.
			\end{enumerate}
		\end{block}
		\begin{exampleblock}{Answer: B}\end{exampleblock}
		\begin{block}{Explanation}
			Portfolio is strategic; inventory is operational. Options A, C, and D misstate the distinction.
		\end{block}
	\end{frame}
	
	\begin{frame}{Q558 --- AI Governance Board Reporting}
		\begin{block}{Question}
			Which is NOT typically included in board-level AI reporting?
			\begin{enumerate}
				\item Aggregate AI risk exposure, including concentration indices, interdependency graphs, and enterprise risk-appetite utilisation across business units.
				\item Fairness and bias metrics, including demographic parity ratios, equalised-odds differentials, and disparate-impact analyses stratified by protected class.
				\item Detailed source code of every model, including line-by-line comments, hyperparameter configurations, gradient histograms, and raw training-data samples.
				\item Material incidents and remediation plans, including root-cause analyses, regulatory notifications, and remediation timelines with accountable owners.
			\end{enumerate}
		\end{block}
		\begin{exampleblock}{Answer: C}\end{exampleblock}
		\begin{block}{Explanation}
			Board reports are strategic and aggregate; source code is not appropriate. Options A, B, and D are typically included.
		\end{block}
	\end{frame}
	
	
	\begin{frame}{Q560 --- AI Risk Taxonomy Categories}
		\begin{block}{Question}
			Which is a common category in AI risk taxonomies?
			\begin{enumerate}
				\item Technical risk, encompassing adversarial robustness, distributional shift, data-provenance integrity, and reproducibility variance.
				\item Ethical risk, encompassing allocative harm, representational harm, transparency deficits, and human-oversight sufficiency.
				\item Operational risk, encompassing third-party concentration, model-drift telemetry, process-failure modes, and resource-exhaustion scenarios.
				\item All of the above are common categories in AI risk taxonomies, alongside legal, regulatory, and reputational dimensions.
			\end{enumerate}
		\end{block}
		\begin{exampleblock}{Answer: D}\end{exampleblock}
		\begin{block}{Explanation}
			AI risk taxonomies span technical, ethical, operational, and other categories; all are common.
		\end{block}
	\end{frame}
	
	% =================================================================
	\section{AI Model Risk in Financial Services}
	% =================================================================
	
	\begin{frame}{Q561 --- Credit Scoring Model Risk}
		\begin{block}{Question}
			Which is the most relevant model risk for an AI-based credit scoring system?
			\begin{enumerate}
				\item The colour scheme and layout of the user interface may not align with brand guidelines, leading to inconsistent customer experience and reduced engagement metrics.
				\item Potential unfairness, inaccuracy, and adverse regulatory consequences from biased or poorly performing predictions, including proxy discrimination, explainability gaps, and ECOA/FCRA non-compliance.
				\item The storage cost of retaining historical training data may exceed budgeted cloud infrastructure allocations over multi-year horizons, impacting FinOps unit economics.
				\item Marketing effectiveness may decline if the model's output is not integrated with campaign management platforms, customer segmentation tools, and CRM workflows.
			\end{enumerate}
		\end{block}
		\begin{exampleblock}{Answer: B}\end{exampleblock}
		\begin{block}{Explanation}
			Credit scoring model risk includes fairness, accuracy, and regulatory risk. Options A, C, and D are secondary or irrelevant.
		\end{block}
	\end{frame}
	
	\begin{frame}{Q562 --- Fraud Detection Model Risk}
		\begin{block}{Question}
			For an AI fraud detection model, which is the most critical model risk?
			\begin{enumerate}
				\item Overfitting to historical fraud patterns, missing novel fraud schemes, and failing to adapt to adversarial concept drift under non-stationary distributions.
				\item Storage cost of retaining transaction-level features and model artefacts across multiple retraining cycles, impacting data-lake tiering strategies.
				\item User interface design that may not clearly present fraud alerts to analysts, slowing investigation workflows and increasing mean-time-to-resolution.
				\item Training time may exceed operational windows, delaying model refresh cycles and reducing responsiveness to emerging fraud typologies.
			\end{enumerate}
		\end{block}
		\begin{exampleblock}{Answer: A}\end{exampleblock}
		\begin{block}{Explanation}
			Overfitting to past patterns can miss novel fraud. Options B, C, and D are less critical.
		\end{block}
	\end{frame}
	
	\begin{frame}{Q563 --- AML Model Risk}
		\begin{block}{Question}
			Which is a key risk for an AI anti-money-laundering (AML) model?
			\begin{enumerate}
				\item Low alert volume may indicate under-detection, reducing investigative workload but increasing regulatory exposure under BSA/AML regimes.
				\item High false-positive rate overwhelming compliance teams and potential missed true positives, both creating regulatory and operational risk.
				\item Storage cost of retaining alert histories and case management records across multiple jurisdictions, impacting data-retention compliance.
				\item Marketing effectiveness may decline if AML alerts are not integrated with customer relationship management platforms and case-management systems.
			\end{enumerate}
		\end{block}
		\begin{exampleblock}{Answer: B}\end{exampleblock}
		\begin{block}{Explanation}
			AML risk includes excessive false positives and missed true positives. Options A, C, and D are secondary.
		\end{block}
	\end{frame}
	
	\begin{frame}{Q564 --- Market Risk Model}
		\begin{block}{Question}
			Which is a typical risk for an AI-enhanced market risk model?
			\begin{enumerate}
				\item Training on non-stationary data leading to unstable predictions and regime-change failures under structural breaks.
				\item Data storage costs may escalate if tick-level market data is retained for extended backtesting periods, impacting data-lake economics.
				\item Logo design may not comply with brand guidelines, creating inconsistent visual identity across risk reports and dashboards.
				\item Employee satisfaction may decline if model outputs are not integrated with existing workflow tools and trading-desk platforms.
			\end{enumerate}
		\end{block}
		\begin{exampleblock}{Answer: A}\end{exampleblock}
		\begin{block}{Explanation}
			Non-stationarity of financial markets is a key challenge. Options B, C, and D are unrelated.
		\end{block}
	\end{frame}
	
	\begin{frame}{Q565 --- Insurance Underwriting AI}
		\begin{block}{Question}
			Which is a specific concern for AI underwriting in insurance?
			\begin{enumerate}
				\item Ethical and legal issues from using alternative data sources that may be considered sensitive or discriminatory under unfair-discrimination statutes.
				\item Faster approvals may reduce customer wait times but could increase adverse selection if not properly controlled via risk-based pricing.
				\item Lower costs may improve competitiveness but could reduce investment in human oversight and quality assurance.
				\item Better customer experience may increase retention but could mask underlying model deficiencies and bias.
			\end{enumerate}
		\end{block}
		\begin{exampleblock}{Answer: A}\end{exampleblock}
		\begin{block}{Explanation}
			Alternative data raises ethical and legal concerns. Options B, C, and D are benefits, not concerns.
		\end{block}
	\end{frame}
	
	\begin{frame}{Q566 --- Model Fairness in Lending}
		\begin{block}{Question}
			Which is a recognised fairness metric used in lending AI?
			\begin{enumerate}
				\item Demographic parity, measured as the ratio of positive prediction rates between protected and non-protected groups.
				\item Storage cost parity, measured as the ratio of data storage costs between protected and non-protected groups.
				\item Latency parity, measured as the ratio of inference times between protected and non-protected groups.
				\item Colour parity, measured as the ratio of UI colour contrast ratios between protected and non-protected groups.
			\end{enumerate}
		\end{block}
		\begin{exampleblock}{Answer: A}\end{exampleblock}
		\begin{block}{Explanation}
			Demographic parity is a standard fairness metric. Options B, C, and D are not.
		\end{block}
	\end{frame}
	
	\begin{frame}{Q567 --- Algorithmic Trading Model Risk}
		\begin{block}{Question}
			Which is a key risk of algorithmic trading models?
			\begin{enumerate}
				\item Flash-crash amplification, unintended feedback loops, and market impact from correlated strategies under stressed liquidity conditions.
				\item Storage capacity may be insufficient to retain tick-level data for regulatory reporting and backtesting, impacting MiFID II compliance.
				\item Logo design may not comply with exchange branding requirements, creating regulatory filing issues and venue-certification delays.
				\item Employee benefits may not be competitive, leading to talent attrition in quantitative teams and loss of institutional knowledge.
			\end{enumerate}
		\end{block}
		\begin{exampleblock}{Answer: A}\end{exampleblock}
		\begin{block}{Explanation}
			Algorithmic trading risks include flash crashes and correlated behaviour. Options B, C, and D are unrelated.
		\end{block}
	\end{frame}
	
	\begin{frame}{Q568 --- Chatbot Risk in Banking}
		\begin{block}{Question}
			Which is the most significant risk from deploying a GenAI chatbot in banking?
			\begin{enumerate}
				\item Generating incorrect or non-compliant advice to customers, creating legal and reputational risk under consumer-protection regimes.
				\item Faster response times may reduce customer wait times but could increase volume of unverified interactions and hallucination exposure.
				\item Cost savings may improve margins but could reduce investment in human oversight and quality assurance.
				\item Customer satisfaction may increase but could mask underlying model deficiencies and bias.
			\end{enumerate}
		\end{block}
		\begin{exampleblock}{Answer: A}\end{exampleblock}
		\begin{block}{Explanation}
			Advice quality and compliance are central concerns. Options B, C, and D are benefits.
		\end{block}
	\end{frame}
	
	\begin{frame}{Q569 --- Model Risk from Alternative Data}
		\begin{block}{Question}
			Which is a key governance concern when using alternative data in financial AI?
			\begin{enumerate}
				\item Legal right to use the data, privacy, and potential for discriminatory impact through proxy variables and disparate-impact mechanisms.
				\item Storage cost may exceed budgeted cloud infrastructure allocations over multi-year horizons, impacting FinOps unit economics.
				\item Model speed may decline if alternative data requires extensive preprocessing and feature engineering pipelines.
				\item Logo design may not comply with brand guidelines, creating inconsistent visual identity across risk reports.
			\end{enumerate}
		\end{block}
		\begin{exampleblock}{Answer: A}\end{exampleblock}
		\begin{block}{Explanation}
			Legal rights, privacy, and fairness are key governance concerns. Options B, C, and D are secondary.
		\end{block}
	\end{frame}
	
	\begin{frame}{Q570 --- Model Risk in AML/KYC}
		\begin{block}{Question}
			What is a specific AI risk in AML/KYC?
			\begin{enumerate}
				\item Customer friction from over-flagging and regulatory exposure from under-flagging, requiring careful threshold calibration under BSA/AML regimes.
				\item Storage cost of retaining KYC documents and alert histories across multiple jurisdictions, impacting data-retention compliance.
				\item Logo design may not comply with brand guidelines, creating inconsistent visual identity across compliance reports.
				\item Training time may exceed operational windows, delaying model refresh cycles and reducing responsiveness to emerging typologies.
			\end{enumerate}
		\end{block}
		\begin{exampleblock}{Answer: A}\end{exampleblock}
		\begin{block}{Explanation}
			Threshold calibration affects both customer experience and regulatory compliance. Options B, C, and D are less critical.
		\end{block}
	\end{frame}
	
	% =================================================================
	\section{AI Model Lifecycle and Change Management}
	% =================================================================
	
	\begin{frame}{Q571 --- Model Lifecycle Stages}
		\begin{block}{Question}
			Which is the correct order of model lifecycle stages?
			\begin{enumerate}
				\item Development, validation, deployment, monitoring, retirement.
				\item Deployment, development, monitoring, validation, retirement.
				\item Monitoring, development, retirement, validation, deployment.
				\item Retirement, development, deployment, monitoring, validation.
			\end{enumerate}
		\end{block}
		\begin{exampleblock}{Answer: A}\end{exampleblock}
		\begin{block}{Explanation}
			Development $\to$ validation $\to$ deployment $\to$ monitoring $\to$ retirement. Options B, C, and D reorder incorrectly.
		\end{block}
	\end{frame}
	
	\begin{frame}{Q572 --- Validation Before Deployment}
		\begin{block}{Question}
			Why must validation occur before deployment?
			\begin{enumerate}
				\item Regulatory convenience, as validation timelines are aligned with regulatory reporting cycles and supervisory review schedules.
				\item To independently confirm the model meets requirements, performs as expected, and is fit for purpose before being used in production.
				\item To reduce cost, as validation is cheaper than post-deployment remediation and incident response.
				\item To increase speed, as validation accelerates deployment timelines through early defect detection.
			\end{enumerate}
		\end{block}
		\begin{exampleblock}{Answer: B}\end{exampleblock}
		\begin{block}{Explanation}
			Pre-deployment validation confirms fitness. Options A, C, and D are secondary.
		\end{block}
	\end{frame}
	
	\begin{frame}{Q573 --- Post-Deployment Monitoring}
		\begin{block}{Question}
			Which activity is part of post-deployment monitoring?
			\begin{enumerate}
				\item Continuous tracking of model performance, data drift, and input integrity, with alerts for threshold breaches and automated escalation.
				\item Only training, as monitoring is limited to retraining cycles and periodic model refreshes.
				\item Only validation, as monitoring is limited to annual validation reviews and independent assessment.
				\item Only code review, as monitoring is limited to static analysis and vulnerability scanning.
			\end{enumerate}
		\end{block}
		\begin{exampleblock}{Answer: A}\end{exampleblock}
		\begin{block}{Explanation}
			Post-deployment monitoring tracks performance and drift. Options B, C, and D are pre-deployment.
		\end{block}
	\end{frame}
	
	\begin{frame}{Q574 --- Recalibration vs Redevelopment}
		\begin{block}{Question}
			Recalibration is preferred to redevelopment when:
			\begin{enumerate}
				\item The model structure is wrong, requiring fundamental redesign and re-engineering of feature pipelines.
				\item The model structure remains valid but parameters need updating using new data under stable distributional assumptions.
				\item The model is retired, as recalibration is only applicable to post-retirement archival models.
				\item The model has never been validated, as recalibration replaces validation and independent review.
			\end{enumerate}
		\end{block}
		\begin{exampleblock}{Answer: B}\end{exampleblock}
		\begin{block}{Explanation}
			Recalibration updates parameters; structural issues require redevelopment. Options A, C, and D are different situations.
		\end{block}
	\end{frame}
	
	\begin{frame}{Q575 --- Model Deprecation}
		\begin{block}{Question}
			What is model deprecation?
			\begin{enumerate}
				\item Immediate retirement, as deprecation is synonymous with decommissioning and archival.
				\item The formal marking of a model as no longer recommended for use, typically while a replacement is being prepared.
				\item Model creation, as deprecation is the initial stage of model development and experimentation.
				\item Model validation, as deprecation is a validation outcome and independent review artefact.
			\end{enumerate}
		\end{block}
		\begin{exampleblock}{Answer: B}\end{exampleblock}
		\begin{block}{Explanation}
			Deprecation flags models for phase-out. Options A, C, and D describe different actions.
		\end{block}
	\end{frame}
	
	\begin{frame}{Q576 --- Model Change Impact Assessment}
		\begin{block}{Question}
			What does a model change impact assessment evaluate?
			\begin{enumerate}
				\item Only cost, as impact assessment is limited to financial considerations and FinOps unit economics.
				\item How a proposed change affects model outputs, downstream users, risk profile, and regulatory compliance.
				\item Only speed, as impact assessment is limited to latency considerations and inference-time optimisation.
				\item Only storage, as impact assessment is limited to infrastructure considerations and data-lake tiering.
			\end{enumerate}
		\end{block}
		\begin{exampleblock}{Answer: B}\end{exampleblock}
		\begin{block}{Explanation}
			Impact assessment covers outputs, downstream, risk, and compliance. Options A, C, and D are too narrow.
		\end{block}
	\end{frame}
	
	\begin{frame}{Q577 --- Model Versioning}
		\begin{block}{Question}
			Why is model versioning important?
			\begin{enumerate}
				\item To reduce storage, as versioning eliminates redundant model artefacts and deduplicates parameter checkpoints.
				\item To track distinct model iterations, enabling reproducibility, audit, and controlled rollback if needed.
				\item To reduce cost, as versioning eliminates redundant training runs and optimises compute allocation.
				\item To speed up training, as versioning accelerates model refresh cycles and hyperparameter search.
			\end{enumerate}
		\end{block}
		\begin{exampleblock}{Answer: B}\end{exampleblock}
		\begin{block}{Explanation}
			Versioning supports reproducibility, audit, and rollback. Options A, C, and D are secondary.
		\end{block}
	\end{frame}
	
	\begin{frame}{Q578 --- Champion/Challenger Framework}
		\begin{block}{Question}
			What is the champion/challenger framework?
			\begin{enumerate}
				\item A marketing strategy, as champion/challenger is used for brand positioning and customer-segment targeting.
				\item A framework where the current production model (champion) is compared against alternative models (challengers) to evaluate whether a switch would improve performance.
				\item A training technique, as champion/challenger is used for hyperparameter tuning and neural architecture search.
				\item A validation report, as champion/challenger is a documentation artefact and independent review deliverable.
			\end{enumerate}
		\end{block}
		\begin{exampleblock}{Answer: B}\end{exampleblock}
		\begin{block}{Explanation}
			Champion/challenger compares production models to alternatives. Options A, C, and D are unrelated.
		\end{block}
	\end{frame}
	
	\begin{frame}{Q579 --- Shadow Deployment}
		\begin{block}{Question}
			What is shadow deployment of a model?
			\begin{enumerate}
				\item Deploying without monitoring, as shadow deployment bypasses monitoring requirements and observability instrumentation.
				\item Running a new model alongside production without using its outputs, to evaluate performance in real conditions.
				\item Deploying to production only, as shadow deployment is a production deployment mode with canary routing.
				\item Retiring a model, as shadow deployment is a decommissioning technique and archival strategy.
			\end{enumerate}
		\end{block}
		\begin{exampleblock}{Answer: B}\end{exampleblock}
		\begin{block}{Explanation}
			Shadow deployment evaluates models in parallel without affecting production. Options A, C, and D misstate.
		\end{block}
	\end{frame}
	
	\begin{frame}{Q580 --- A/B Testing of Models}
		\begin{block}{Question}
			What is A/B testing in the AI deployment context?
			\begin{enumerate}
				\item Randomly assigning users to two model versions and comparing outcomes under controlled experimental conditions.
				\item Retiring both models, as A/B testing is a decommissioning technique and archival strategy.
				\item Doubling training data, as A/B testing is a data augmentation technique and synthetic-data generation method.
				\item Storing data twice, as A/B testing is a storage redundancy technique and disaster-recovery strategy.
			\end{enumerate}
		\end{block}
		\begin{exampleblock}{Answer: A}\end{exampleblock}
		\begin{block}{Explanation}
			A/B testing compares model versions via controlled experiments. Options B, C, and D are unrelated.
		\end{block}
	\end{frame}
	
	% =================================================================
	\section{AI Incident Management}
	% =================================================================
	
	\begin{frame}{Q581 --- AI Incident Response}
		\begin{block}{Question}
			Which is a key element of an AI incident response plan?
			\begin{enumerate}
				\item Ignoring incidents, as incident response is limited to post-mortem documentation and retrospective analysis.
				\item Defined escalation paths, roles, playbooks, and communication plans for handling AI-related incidents.
				\item Only post-mortems, as incident response is limited to retrospective analysis and lessons-learned exercises.
				\item Only IT involvement, as incident response is limited to infrastructure teams and site-reliability engineers.
			\end{enumerate}
		\end{block}
		\begin{exampleblock}{Answer: B}\end{exampleblock}
		\begin{block}{Explanation}
			An incident response plan includes escalation, roles, playbooks, and communication. Options A, C, and D are insufficient.
		\end{block}
	\end{frame}
	
	\begin{frame}{Q582 --- AI Incident Root Cause}
		\begin{block}{Question}
			An AI model begins producing biased outputs. Which is the most likely root cause?
			\begin{enumerate}
				\item Data drift or newly biased data, as the model's input distribution has shifted under non-stationary conditions.
				\item Server failure, as hardware faults can introduce bias into model outputs through corrupted tensor operations.
				\item Power outage, as interrupted training can introduce bias into model weights through checkpoint corruption.
				\item UI redesign, as changes to the user interface can introduce bias into model inputs through altered data-capture flows.
			\end{enumerate}
		\end{block}
		\begin{exampleblock}{Answer: A}\end{exampleblock}
		\begin{block}{Explanation}
			Bias in outputs typically stems from data or model development. Options B, C, and D are unrelated.
		\end{block}
	\end{frame}
	
	\begin{frame}{Q583 --- Post-Incident Review}
		\begin{block}{Question}
			What is the purpose of a post-incident review for AI incidents?
			\begin{enumerate}
				\item Assign blame, as post-incident reviews are primarily disciplinary processes and accountability mechanisms.
				\item Learn from the incident, improve controls and processes, and prevent recurrence.
				\item End investigation, as post-incident reviews are terminal documentation exercises and closure artefacts.
				\item Skip documentation, as post-incident reviews are informal debriefs and verbal knowledge-transfer sessions.
			\end{enumerate}
		\end{block}
		\begin{exampleblock}{Answer: B}\end{exampleblock}
		\begin{block}{Explanation}
			Post-incident reviews focus on learning and prevention. Options A, C, and D are contrary.
		\end{block}
	\end{frame}
	
	\begin{frame}{Q584 --- Incident Communication}
		\begin{block}{Question}
			Why is communication important during an AI incident?
			\begin{enumerate}
				\item It is not important, as incident response is purely technical and does not require stakeholder engagement.
				\item Clear communication with regulators, customers, and stakeholders maintains trust and meets regulatory obligations.
				\item Only internal communication matters, as external parties are not affected by AI incidents.
				\item Only external communication matters, as internal teams are not affected by AI incidents.
			\end{enumerate}
		\end{block}
		\begin{exampleblock}{Answer: B}\end{exampleblock}
		\begin{block}{Explanation}
			Communication with all stakeholders is essential. Options A, C, and D are incomplete or wrong.
		\end{block}
	\end{frame}
	
	\begin{frame}{Q585 --- AI Incident Classification}
		\begin{block}{Question}
			How are AI incidents typically classified?
			\begin{enumerate}
				\item By colour, as incident classification uses a colour-coded severity scale and RAG status indicators.
				\item By severity and impact, considering factors such as harm to customers, financial loss, and regulatory exposure.
				\item Randomly, as incident classification is not standardised and varies by team discretion.
				\item By alphabet, as incident classification uses alphabetical ordering and lexical sorting.
			\end{enumerate}
		\end{block}
		\begin{exampleblock}{Answer: B}\end{exampleblock}
		\begin{block}{Explanation}
			Incidents are classified by severity and impact. Options A, C, and D are not standard.
		\end{block}
	\end{frame}
	
	\begin{frame}{Q586 --- Regulatory Notification}
		\begin{block}{Question}
			When a serious AI incident affects customers, what is typically required?
			\begin{enumerate}
				\item No action, as AI incidents are not subject to regulatory notification and supervisory reporting.
				\item Prompt notification to regulators where required and remediation plans.
				\item Only internal review, as regulatory notification is optional and discretionary.
				\item Silence, as public disclosure is prohibited and confidentiality is paramount.
			\end{enumerate}
		\end{block}
		\begin{exampleblock}{Answer: B}\end{exampleblock}
		\begin{block}{Explanation}
			Serious incidents may trigger regulatory notification. Options A, C, and D violate governance.
		\end{block}
	\end{frame}
	
	\begin{frame}{Q587 --- Incident Metrics}
		\begin{block}{Question}
			Which metric helps track AI incident management effectiveness?
			\begin{enumerate}
				\item Total storage used, as storage metrics correlate with incident management effectiveness and data-retention efficiency.
				\item Time to detection, time to remediation, recurrence rate, and impact.
				\item Marketing reach, as customer awareness correlates with incident management effectiveness and brand perception.
				\item UI clicks, as user engagement correlates with incident management effectiveness and interface usability.
			\end{enumerate}
		\end{block}
		\begin{exampleblock}{Answer: B}\end{exampleblock}
		\begin{block}{Explanation}
			Detection, remediation, recurrence, and impact are key metrics. Options A, C, and D are unrelated.
		\end{block}
	\end{frame}
	
	\begin{frame}{Q588 --- Red Teaming}
		\begin{block}{Question}
			What is red teaming in AI?
			\begin{enumerate}
				\item A marketing exercise, as red teaming is used for brand positioning and competitive differentiation.
				\item Adversarial testing by a dedicated team to uncover vulnerabilities and unintended behaviours before deployment.
				\item A financial audit, as red teaming is used for financial controls testing and SOX compliance.
				\item A user survey, as red teaming is used for customer feedback collection and satisfaction measurement.
			\end{enumerate}
		\end{block}
		\begin{exampleblock}{Answer: B}\end{exampleblock}
		\begin{block}{Explanation}
			Red teaming proactively tests for vulnerabilities. Options A, C, and D misstate.
		\end{block}
	\end{frame}
	
	\begin{frame}{Q589 --- AI Incident Data}
		\begin{block}{Question}
			Why is logging data important for AI incident management?
			\begin{enumerate}
				\item To increase storage, as logging is primarily a storage optimization technique and data-lake tiering strategy.
				\item To reconstruct events, identify root causes, and support regulatory reporting.
				\item To slow down systems, as logging introduces latency for diagnostic purposes and observability instrumentation.
				\item To increase costs, as logging is primarily a cost optimization technique and FinOps unit-economics lever.
			\end{enumerate}
		\end{block}
		\begin{exampleblock}{Answer: B}\end{exampleblock}
		\begin{block}{Explanation}
			Logging supports root-cause analysis and reporting. Options A, C, and D are incorrect or secondary.
		\end{block}
	\end{frame}
	
	\begin{frame}{Q590 --- Continuous Improvement}
		\begin{block}{Question}
			What is the purpose of continuous improvement in AI governance?
			\begin{enumerate}
				\item To increase costs, as continuous improvement is primarily a cost optimization technique and FinOps unit-economics lever.
				\item To iterate on governance frameworks and controls based on incidents, audits, and evolving best practices.
				\item To avoid change, as continuous improvement is primarily a stability technique and change-management control.
				\item To slow down deployment, as continuous improvement is primarily a risk mitigation technique and deployment-gating mechanism.
			\end{enumerate}
		\end{block}
		\begin{exampleblock}{Answer: B}\end{exampleblock}
		\begin{block}{Explanation}
			Continuous improvement ensures governance keeps pace with change. Options A, C, and D are contrary.
		\end{block}
	\end{frame}
	
	% =================================================================
	\section{AI Governance Maturity}
	% =================================================================
	
	\begin{frame}{Q591 --- AI Governance Maturity Model}
		\begin{block}{Question}
			A typical AI governance maturity model includes levels such as:
			\begin{enumerate}
				\item Ad hoc, developing, defined, managed, optimised.
				\item Red, blue, green, as maturity models use colour-coded levels and RAG status indicators.
				\item Only one level, as maturity models are binary and non-hierarchical.
				\item No levels, as maturity models are non-hierarchical and network-based.
			\end{enumerate}
		\end{block}
		\begin{exampleblock}{Answer: A}\end{exampleblock}
		\begin{block}{Explanation}
			Maturity models typically progress from ad hoc to optimised. Options B, C, and D are not standard.
		\end{block}
	\end{frame}
	
	\begin{frame}{Q592 --- AI Governance Gaps}
		\begin{block}{Question}
			Which is a common gap in AI governance frameworks?
			\begin{enumerate}
				\item Extensive controls, as frameworks often over-control low-risk use cases and create unnecessary friction.
				\item Missing governance for GenAI, alternative data, or third-party models.
				\item Complete coverage, as frameworks often over-cover all risk dimensions and create unnecessary complexity.
				\item Perfect alignment, as frameworks often over-align with business objectives and create unnecessary constraints.
			\end{enumerate}
		\end{block}
		\begin{exampleblock}{Answer: B}\end{exampleblock}
		\begin{block}{Explanation}
			Common gaps include GenAI, alternative data, and vendor models. Options A, C, and D are not gaps.
		\end{block}
	\end{frame}
	
	\begin{frame}{Q593 --- Governance Roadmap}
		\begin{block}{Question}
			What is the purpose of an AI governance roadmap?
			\begin{enumerate}
				\item A marketing document, as roadmaps are used for brand positioning and competitive differentiation.
				\item A phased plan outlining how the organisation will build, enhance, and sustain AI governance capabilities over time.
				\item A single policy, as roadmaps are policy documents and compliance artefacts.
				\item A code repository, as roadmaps are technical artefacts and infrastructure-as-code modules.
			\end{enumerate}
		\end{block}
		\begin{exampleblock}{Answer: B}\end{exampleblock}
		\begin{block}{Explanation}
			A roadmap is a phased plan. Options A, C, and D misstate.
		\end{block}
	\end{frame}
	
	\begin{frame}{Q594 --- AI Governance KPIs}
		\begin{block}{Question}
			Which is an appropriate KPI for AI governance?
			\begin{enumerate}
				\item Number of employees, as headcount correlates with governance effectiveness and organisational capacity.
				\item Percentage of models validated on schedule.
				\item Office size, as physical footprint correlates with governance effectiveness and operational scale.
				\item Marketing spend, as brand awareness correlates with governance effectiveness and market positioning.
			\end{enumerate}
		\end{block}
		\begin{exampleblock}{Answer: B}\end{exampleblock}
		\begin{block}{Explanation}
			Validation timeliness is a governance KPI. Options A, C, and D are unrelated.
		\end{block}
	\end{frame}
	
	\begin{frame}{Q595 --- AI Governance Training}
		\begin{block}{Question}
			Which audience requires AI governance training?
			\begin{enumerate}
				\item Only data scientists, as governance training is limited to technical roles and model-development functions.
				\item A broad set including executives, developers, validators, business users, compliance, and internal audit.
				\item Only IT, as governance training is limited to infrastructure teams and site-reliability engineers.
				\item Only external auditors, as governance training is limited to assurance providers and independent reviewers.
			\end{enumerate}
		\end{block}
		\begin{exampleblock}{Answer: B}\end{exampleblock}
		\begin{block}{Explanation}
			AI governance training spans many roles. Options A, C, and D are too narrow.
		\end{block}
	\end{frame}
	
	\begin{frame}{Q596 --- Governance Data Collection}
		\begin{block}{Question}
			Why is data collection for AI governance metrics important?
			\begin{enumerate}
				\item To meet marketing goals, as governance metrics are used for brand positioning and competitive differentiation.
				\item To enable measurement, benchmarking, and continuous improvement of AI governance practices.
				\item To increase costs, as governance metrics are primarily a cost optimization technique and FinOps unit-economics lever.
				\item To slow operations, as governance metrics are primarily a risk mitigation technique and deployment-gating mechanism.
			\end{enumerate}
		\end{block}
		\begin{exampleblock}{Answer: B}\end{exampleblock}
		\begin{block}{Explanation}
			Data supports measurement and improvement. Options A, C, and D are incorrect.
		\end{block}
	\end{frame}
	
	\begin{frame}{Q597 --- AI Governance Certification}
		\begin{block}{Question}
			What is the purpose of AI governance certifications (e.g., ISO/IEC 42001)?
			\begin{enumerate}
				\item Marketing only, as certifications are used for brand positioning and competitive differentiation.
				\item To demonstrate externally that the organisation meets recognised standards for AI governance.
				\item To replace regulations, as certifications are legally binding substitutes and safe-harbour mechanisms.
				\item To lower costs, as certifications are primarily a cost optimization technique and FinOps unit-economics lever.
			\end{enumerate}
		\end{block}
		\begin{exampleblock}{Answer: B}\end{exampleblock}
		\begin{block}{Explanation}
			Certifications demonstrate conformance with standards. Options A, C, and D are not the primary purpose.
		\end{block}
	\end{frame}
	
	\begin{frame}{Q598 --- Board AI Literacy}
		\begin{block}{Question}
			Why is board AI literacy important?
			\begin{enumerate}
				\item To increase headcount, as board literacy correlates with hiring needs and organisational capacity.
				\item To enable effective oversight of AI strategy, risk, and governance at the highest level.
				\item To reduce costs, as board literacy correlates with cost efficiency and FinOps unit economics.
				\item To speed deployment, as board literacy correlates with deployment velocity and time-to-market.
			\end{enumerate}
		\end{block}
		\begin{exampleblock}{Answer: B}\end{exampleblock}
		\begin{block}{Explanation}
			Board AI literacy enables effective oversight. Options A, C, and D are not the primary reason.
		\end{block}
	\end{frame}
	
	\begin{frame}{Q599 --- AI Ethics Committee Integration}
		\begin{block}{Question}
			How should an AI ethics committee integrate with the model risk function?
			\begin{enumerate}
				\item No integration is needed, as ethics and model risk are independent domains and separate governance tracks.
				\item Through coordinated review of high-risk AI use cases, sharing of insights on fairness and ethics, and joint escalation to senior management.
				\item Only through annual meetings, as integration is limited to periodic reviews and annual attestations.
				\item Only through external advisors, as integration is outsourced to third parties and independent consultants.
			\end{enumerate}
		\end{block}
		\begin{exampleblock}{Answer: B}\end{exampleblock}
		\begin{block}{Explanation}
			Coordination between ethics and model risk is essential. Options A, C, and D are insufficient.
		\end{block}
	\end{frame}
	
	\begin{frame}{Q600 --- Governance Evolution}
		\begin{block}{Question}
			Why must AI governance frameworks evolve?
			\begin{enumerate}
				\item To increase cost, as evolution is primarily a cost optimization technique and FinOps unit-economics lever.
				\item Because AI technologies, use cases, and regulations change rapidly, requiring ongoing updates to remain effective.
				\item To slow operations, as evolution is primarily a risk mitigation technique and deployment-gating mechanism.
				\item To reduce transparency, as evolution is primarily a confidentiality technique and information-security control.
			\end{enumerate}
		\end{block}
		\begin{exampleblock}{Answer: B}\end{exampleblock}
		\begin{block}{Explanation}
			Rapid change in AI and regulation requires evolving governance. Options A, C, and D contradict this.
		\end{block}
	\end{frame}
	
	% =================================================================
	\section{Ethics in GenAI}
	% =================================================================
	
	\begin{frame}{Q601 --- GenAI and Autonomy}
		\begin{block}{Question}
			How can GenAI undermine human autonomy?
			\begin{enumerate}
				\item By providing choices, as GenAI expands the option set available to users and increases decision-space dimensionality.
				\item By generating persuasive, tailored content that manipulates users without their awareness.
				\item By explaining decisions, as GenAI provides transparency into decision processes and interpretability artefacts.
				\item By slowing systems, as GenAI introduces latency that reduces user control and increases interaction friction.
			\end{enumerate}
		\end{block}
		\begin{exampleblock}{Answer: B}\end{exampleblock}
		\begin{block}{Explanation}
			Manipulative content undermines autonomy. Options A, C, and D respect or are unrelated to autonomy.
		\end{block}
	\end{frame}
	
	\begin{frame}{Q602 --- GenAI and Bias}
		\begin{block}{Question}
			How does GenAI amplify bias?
			\begin{enumerate}
				\item By generating content that reflects biases present in training data, at scale.
				\item By reducing accuracy, as lower accuracy correlates with higher bias and increased error rates.
				\item By slowing training, as longer training correlates with higher bias and overfitting to minority classes.
				\item By increasing costs, as higher costs correlate with higher bias and increased computational complexity.
			\end{enumerate}
		\end{block}
		\begin{exampleblock}{Answer: A}\end{exampleblock}
		\begin{block}{Explanation}
			GenAI can amplify training-data bias at scale. Options B, C, and D are unrelated.
		\end{block}
	\end{frame}
	
	\begin{frame}{Q603 --- GenAI and Truth}
		\begin{block}{Question}
			Why is GenAI a challenge for truth and accountability?
			\begin{enumerate}
				\item It only outputs truth, as GenAI is designed for factual accuracy and grounded generation.
				\item It generates plausible but sometimes incorrect content, making accountability for misinformation difficult.
				\item It has no training data, as GenAI is designed for zero-shot learning and few-shot generalisation.
				\item It never hallucinates, as GenAI is designed for factual accuracy and retrieval-augmented grounding.
			\end{enumerate}
		\end{block}
		\begin{exampleblock}{Answer: B}\end{exampleblock}
		\begin{block}{Explanation}
			Hallucinations and plausible output complicate accountability. Options A, C, and D are incorrect.
		\end{block}
	\end{frame}
	
	\begin{frame}{Q604 --- GenAI and Consent}
		\begin{block}{Question}
			What consent issues arise with GenAI training data?
			\begin{enumerate}
				\item None, as GenAI training data is always consented and properly licensed.
				\item Data subjects whose content was used in training may not have consented, raising ethical and legal concerns.
				\item Consent is automatically granted, as GenAI training data is publicly available and web-scraped.
				\item Consent is not required, as GenAI training data is exempt from consent requirements under fair-use doctrines.
			\end{enumerate}
		\end{block}
		\begin{exampleblock}{Answer: B}\end{exampleblock}
		\begin{block}{Explanation}
			Training data consent is a critical concern. Options A, C, and D are incorrect.
		\end{block}
	\end{frame}
	
	\begin{frame}{Q605 --- GenAI and Attribution}
		\begin{block}{Question}
			Why is attribution important in GenAI?
			\begin{enumerate}
				\item It isn't important, as GenAI outputs are not attributable and are generated de novo.
				\item To trace sources of content and information, supporting transparency and possibly legal rights such as copyright.
				\item To increase cost, as attribution is primarily a cost optimization technique and FinOps unit-economics lever.
				\item To slow inference, as attribution is primarily a latency optimization technique and inference-time acceleration method.
			\end{enumerate}
		\end{block}
		\begin{exampleblock}{Answer: B}\end{exampleblock}
		\begin{block}{Explanation}
			Attribution supports transparency and copyright compliance. Options A, C, and D are incorrect.
		\end{block}
	\end{frame}
	
	\begin{frame}{Q606 --- GenAI and Copyright}
		\begin{block}{Question}
			What is a copyright concern with GenAI?
			\begin{enumerate}
				\item None, as GenAI outputs are always copyright-clean and free of encumbrances.
				\item Training on copyrighted material without permission and generating content that may infringe copyright.
				\item GenAI is always copyright-clean, as training data is licensed and properly attributed.
				\item Copyright does not apply, as GenAI outputs are not fixed in a tangible medium and are ephemeral.
			\end{enumerate}
		\end{block}
		\begin{exampleblock}{Answer: B}\end{exampleblock}
		\begin{block}{Explanation}
			Copyright concerns are significant. Options A, C, and D are incorrect.
		\end{block}
	\end{frame}
	
	\begin{frame}{Q607 --- GenAI and Safety}
		\begin{block}{Question}
			Which is a safety concern with GenAI in critical applications?
			\begin{enumerate}
				\item Faster response, as speed reduces safety in critical applications and increases error propagation.
				\item Generating unsafe or harmful content that could cause physical, psychological, or financial harm.
				\item Lower costs, as cost reduction reduces safety in critical applications and increases technical debt.
				\item Higher accuracy, as accuracy reduces safety in critical applications and increases overconfidence.
			\end{enumerate}
		\end{block}
		\begin{exampleblock}{Answer: B}\end{exampleblock}
		\begin{block}{Explanation}
			Harmful content generation is a safety concern. Options A, C, and D are benefits.
		\end{block}
	\end{frame}
	
	\begin{frame}{Q608 --- GenAI and Consumer Protection}
		\begin{block}{Question}
			Which consumer protection issue arises from GenAI?
			\begin{enumerate}
				\item Increasing prices only, as GenAI directly causes price increases and inflationary pressure.
				\item Misleading or inaccurate information provided to consumers, potentially leading to harm.
				\item Faster service, as speed reduces consumer protection and increases error propagation.
				\item Lower fees, as fee reduction reduces consumer protection and increases hidden costs.
			\end{enumerate}
		\end{block}
		\begin{exampleblock}{Answer: B}\end{exampleblock}
		\begin{block}{Explanation}
			Misleading information is a consumer protection issue. Options A, C, and D are unrelated.
		\end{block}
	\end{frame}
	
	\begin{frame}{Q609 --- GenAI and Financial Advice}
		\begin{block}{Question}
			What regulatory concern arises when GenAI is used for financial advice?
			\begin{enumerate}
				\item Lower costs, as cost reduction reduces regulatory compliance and increases risk-taking.
				\item Potential non-compliance with advice regulations if unvetted or unsuitable advice is provided.
				\item Faster response, as speed reduces regulatory compliance and increases error propagation.
				\item Higher throughput, as throughput increases regulatory compliance and reduces oversight.
			\end{enumerate}
		\end{block}
		\begin{exampleblock}{Answer: B}\end{exampleblock}
		\begin{block}{Explanation}
			Regulatory concerns exist around advice suitability. Options A, C, and D are benefits, not concerns.
		\end{block}
	\end{frame}
	
	\begin{frame}{Q610 --- GenAI and Fraud}
		\begin{block}{Question}
			How can GenAI facilitate fraud?
			\begin{enumerate}
				\item By making it easier to create convincing phishing content and deepfakes at scale.
				\item By improving accuracy, as accuracy reduces fraud and increases detection rates.
				\item By reducing costs, as cost reduction reduces fraud and increases detection efficiency.
				\item By speeding up transactions, as speed reduces fraud and increases detection velocity.
			\end{enumerate}
		\end{block}
		\begin{exampleblock}{Answer: A}\end{exampleblock}
		\begin{block}{Explanation}
			GenAI enables scalable, convincing fraud content. Options B, C, and D are unrelated.
		\end{block}
	\end{frame}
	
	% =================================================================
	\section{Emerging Topics}
	% =================================================================
	
	\begin{frame}{Q611 --- Agentic AI Risks}
		\begin{block}{Question}
			What are the key risks of agentic AI?
			\begin{enumerate}
				\item Unintended autonomous actions, difficulty in oversight, and cascading effects from chains of decisions.
				\item Faster decisions, as speed reduces oversight requirements and increases error propagation.
				\item Lower cost, as cost reduction reduces oversight requirements and increases technical debt.
				\item Simplicity, as simplicity reduces oversight requirements and increases hidden complexity.
			\end{enumerate}
		\end{block}
		\begin{exampleblock}{Answer: A}\end{exampleblock}
		\begin{block}{Explanation}
			Agentic AI risks relate to autonomy and oversight. Options B, C, and D are benefits or inaccurate.
		\end{block}
	\end{frame}
	
	\begin{frame}{Q612 --- Multi-Agent Systems}
		\begin{block}{Question}
			What risk is specific to multi-agent AI systems?
			\begin{enumerate}
				\item Faster training, as speed reduces coordination risk and increases emergent complexity.
				\item Emergent behaviour arising from interactions between agents, which may be difficult to predict or control.
				\item Lower costs, as cost reduction reduces coordination risk and increases technical debt.
				\item Higher accuracy, as accuracy reduces coordination risk and increases overconfidence.
			\end{enumerate}
		\end{block}
		\begin{exampleblock}{Answer: B}\end{exampleblock}
		\begin{block}{Explanation}
			Emergent multi-agent behaviour is a specific risk. Options A, C, and D are benefits.
		\end{block}
	\end{frame}
	
	\begin{frame}{Q613 --- AI and Cyber Risk}
		\begin{block}{Question}
			How does AI increase cyber risk?
			\begin{enumerate}
				\item By reducing attack surface, as AI simplifies infrastructure and reduces vulnerability exposure.
				\item By enabling more sophisticated, scalable attacks, and by introducing new vulnerabilities (e.g., model extraction, adversarial inputs).
				\item By eliminating threats, as AI replaces traditional security controls and reduces human error.
				\item By lowering costs, as AI reduces security investment needs and increases automation coverage.
			\end{enumerate}
		\end{block}
		\begin{exampleblock}{Answer: B}\end{exampleblock}
		\begin{block}{Explanation}
			AI introduces new attack vectors and scales threats. Options A, C, and D are incorrect.
		\end{block}
	\end{frame}
	
	\begin{frame}{Q614 --- AI Supply Chain Risk}
		\begin{block}{Question}
			What is AI supply chain risk?
			\begin{enumerate}
				\item Risk of low prices, as low-cost suppliers may reduce quality and increase technical debt.
				\item Risk arising from dependencies on external AI providers, datasets, and libraries, including security, availability, and continuity concerns.
				\item Lower costs, as cost reduction reduces supply chain risk and increases vendor concentration.
				\item Higher speed, as speed reduces supply chain risk and increases integration complexity.
			\end{enumerate}
		\end{block}
		\begin{exampleblock}{Answer: B}\end{exampleblock}
		\begin{block}{Explanation}
			Supply chain risk involves external dependencies. Options A, C, and D are benefits or unrelated.
		\end{block}
	\end{frame}
	
	\begin{frame}{Q615 --- AI and Third-Party Risk}
		\begin{block}{Question}
			How should third-party AI risk be managed?
			\begin{enumerate}
				\item Ignore it, as third-party risk is outside the organisation's control and cannot be mitigated.
				\item Through due diligence, contractual protections, ongoing monitoring, and contingency planning for provider failures.
				\item Only contracts, as contractual protections are sufficient and eliminate residual risk.
				\item Only internal audits, as internal audits cover third-party risk and provide independent assurance.
			\end{enumerate}
		\end{block}
		\begin{exampleblock}{Answer: B}\end{exampleblock}
		\begin{block}{Explanation}
			Third-party risk requires due diligence, contracts, monitoring, and contingency planning. Options A, C, and D are insufficient.
		\end{block}
	\end{frame}
	
	\begin{frame}{Q616 --- AI and Concentration Risk}
		\begin{block}{Question}
			Why is AI concentration risk a concern?
			\begin{enumerate}
				\item Because it reduces costs, as concentration reduces procurement overhead and increases bargaining power.
				\item Because reliance on a small number of AI providers or models creates correlated failure risk across the ecosystem.
				\item Because it increases accuracy, as concentration improves model quality and increases benchmark performance.
				\item Because it improves speed, as concentration reduces integration complexity and increases deployment velocity.
			\end{enumerate}
		\end{block}
		\begin{exampleblock}{Answer: B}\end{exampleblock}
		\begin{block}{Explanation}
			Concentration amplifies correlated risk. Options A, C, and D are benefits.
		\end{block}
	\end{frame}
	
	\begin{frame}{Q617 --- AI and Reputational Risk}
		\begin{block}{Question}
			Which is a reputational risk from AI?
			\begin{enumerate}
				\item Faster delivery, as speed reduces reputational risk and increases customer satisfaction.
				\item Public backlash from biased, unethical, or incompetent AI use, resulting in loss of trust and brand value.
				\item Lower cost, as cost reduction reduces reputational risk and increases competitiveness.
				\item Higher accuracy, as accuracy reduces reputational risk and increases customer trust.
			\end{enumerate}
		\end{block}
		\begin{exampleblock}{Answer: B}\end{exampleblock}
		\begin{block}{Explanation}
			Reputational risk stems from public perception of AI failures. Options A, C, and D are benefits.
		\end{block}
	\end{frame}
	
	\begin{frame}{Q618 --- AI and Legal Risk}
		\begin{block}{Question}
			Which is a legal risk from AI?
			\begin{enumerate}
				\item Lower cost, as cost reduction reduces legal risk and increases compliance efficiency.
				\item Liability for harm caused by AI decisions, non-compliance with regulations, and IP violations.
				\item Faster deployment, as speed reduces legal risk and increases time-to-market.
				\item Higher accuracy, as accuracy reduces legal risk and increases regulatory compliance.
			\end{enumerate}
		\end{block}
		\begin{exampleblock}{Answer: B}\end{exampleblock}
		\begin{block}{Explanation}
			Legal risks include liability, non-compliance, and IP issues. Options A, C, and D are benefits.
		\end{block}
	\end{frame}
	
	\begin{frame}{Q619 --- AI and Environmental Risk}
		\begin{block}{Question}
			What environmental risk does AI present?
			\begin{enumerate}
				\item None, as AI has no environmental impact and is carbon-neutral by design.
				\item Significant compute-related energy consumption and carbon emissions from training and inference of large models.
				\item Lower energy use, as AI reduces energy consumption and increases efficiency across sectors.
				\item Positive environmental impact only, as AI enables sustainability and reduces carbon footprints.
			\end{enumerate}
		\end{block}
		\begin{exampleblock}{Answer: B}\end{exampleblock}
		\begin{block}{Explanation}
			AI compute consumption creates environmental risk. Options A, C, and D are incorrect.
		\end{block}
	\end{frame}
	
	\begin{frame}{Q620 --- AI and Regulatory Fragmentation}
		\begin{block}{Question}
			Why is regulatory fragmentation a challenge for AI governance?
			\begin{enumerate}
				\item Because it speeds up compliance, as fragmentation reduces regulatory burden and increases harmonisation.
				\item Because different jurisdictions impose different rules, complicating compliance for globally operating firms.
				\item Because it reduces cost, as fragmentation reduces compliance overhead and increases efficiency.
				\item Because it simplifies everything, as fragmentation reduces complexity and increases clarity.
			\end{enumerate}
		\end{block}
		\begin{exampleblock}{Answer: B}\end{exampleblock}
		\begin{block}{Explanation}
			Fragmented regulation complicates global compliance. Options A, C, and D are incorrect.
		\end{block}
	\end{frame}
	
	% =================================================================
	\section{Integrating AI Governance into the Three Lines}
	% =================================================================
	
	\begin{frame}{Q621 --- First Line AI Responsibilities}
		\begin{block}{Question}
			Which is the first line's responsibility for AI models?
			\begin{enumerate}
				\item Independent validation, as the first line validates models independently and provides objective assessment.
				\item Design, develop, own, and operate AI models; monitor performance and own data quality.
				\item Set policy, as the first line sets AI governance policy and establishes risk appetite.
				\item Audit, as the first line audits AI models and provides independent assurance.
			\end{enumerate}
		\end{block}
		\begin{exampleblock}{Answer: B}\end{exampleblock}
		\begin{block}{Explanation}
			First line owns AI models and operations. Options A, C, and D belong to second or third lines.
		\end{block}
	\end{frame}
	
	\begin{frame}{Q622 --- Second Line AI Responsibilities}
		\begin{block}{Question}
			Which is the second line's responsibility for AI models?
			\begin{enumerate}
				\item Build AI models, as the second line builds models and owns data pipelines.
				\item Independent validation, policy, and challenge of first-line assumptions and controls.
				\item Use models for decisions, as the second line uses models for business decisions and customer interactions.
				\item Market models, as the second line markets AI models and drives adoption.
			\end{enumerate}
		\end{block}
		\begin{exampleblock}{Answer: B}\end{exampleblock}
		\begin{block}{Explanation}
			Second line validates, sets policy, and challenges. Options A, C, and D belong to first line.
		\end{block}
	\end{frame}
	
	\begin{frame}{Q623 --- Third Line AI Responsibilities}
		\begin{block}{Question}
			Which is the third line's responsibility for AI models?
			\begin{enumerate}
				\item Build models, as the third line builds models and owns data pipelines.
				\item Provide independent assurance on the effectiveness of AI governance and controls.
				\item Set risk appetite, as the third line sets risk appetite and establishes tolerances.
				\item Manage model data, as the third line manages model data and owns data quality.
			\end{enumerate}
		\end{block}
		\begin{exampleblock}{Answer: B}\end{exampleblock}
		\begin{block}{Explanation}
			Third line provides independent assurance. Options A, C, and D belong to first or second lines.
		\end{block}
	\end{frame}
	
	\begin{frame}{Q624 --- Independence in AI Validation}
		\begin{block}{Question}
			Why must AI validation be independent?
			\begin{enumerate}
				\item To speed up processes, as independence accelerates validation timelines and reduces bureaucratic friction.
				\item To avoid conflicts of interest and ensure objective assessment of model quality and risk.
				\item To reduce costs, as independence reduces validation costs and increases efficiency.
				\item To bypass policies, as independence exempts validation from policies and standard procedures.
			\end{enumerate}
		\end{block}
		\begin{exampleblock}{Answer: B}\end{exampleblock}
		\begin{block}{Explanation}
			Independence ensures objective validation. Options A, C, and D are incorrect.
		\end{block}
	\end{frame}
	
	\begin{frame}{Q625 --- Effective Challenge in AI}
		\begin{block}{Question}
			What does effective challenge mean in the AI governance context?
			\begin{enumerate}
				\item Automatic approval, as effective challenge is a formality and routine step in the approval workflow.
				\item Critical, objective, and informed review of AI models and their assumptions by qualified parties, independent from developers.
				\item Marketing endorsement, as effective challenge is a promotional activity and brand-positioning exercise.
				\item Rubber-stamping, as effective challenge is a compliance formality and documentation exercise.
			\end{enumerate}
		\end{block}
		\begin{exampleblock}{Answer: B}\end{exampleblock}
		\begin{block}{Explanation}
			Effective challenge is critical review. Options A, C, and D are inconsistent.
		\end{block}
	\end{frame}
	
	\begin{frame}{Q626 --- AI Risk Appetite Cascading}
		\begin{block}{Question}
			How should AI risk appetite cascade through the organisation?
			\begin{enumerate}
				\item It shouldn't, as risk appetite is set independently at each level and is not hierarchical.
				\item From board-set appetite to business-unit tolerances and model-level thresholds, all aligned.
				\item Only at business-unit level, as board appetite is not relevant to AI and is too high-level.
				\item Only at IT level, as AI risk is an IT concern and is managed by technology teams.
			\end{enumerate}
		\end{block}
		\begin{exampleblock}{Answer: B}\end{exampleblock}
		\begin{block}{Explanation}
			Appetite cascades from board to units to models. Options A, C, and D are too narrow.
		\end{block}
	\end{frame}
	
	\begin{frame}{Q627 --- AI Risk Ownership}
		\begin{block}{Question}
			Who owns AI risk in an organisation?
			\begin{enumerate}
				\item The IT department, as AI risk is an IT concern and is managed by technology teams.
				\item Business units that use AI, supported by risk, compliance, and other functions.
				\item Only the CRO, as AI risk is a risk function concern and is managed by risk teams.
				\item Only external vendors, as AI risk is a vendor concern and is managed by third parties.
			\end{enumerate}
		\end{block}
		\begin{exampleblock}{Answer: B}\end{exampleblock}
		\begin{block}{Explanation}
			AI risk ownership sits with business users supported by risk functions. Options A, C, and D are incomplete.
		\end{block}
	\end{frame}
	
	\begin{frame}{Q628 --- AI Escalation Paths}
		\begin{block}{Question}
			What is the purpose of AI escalation paths?
			\begin{enumerate}
				\item To slow down decisions, as escalation introduces delays and bureaucratic friction.
				\item To ensure that AI incidents, breaches, and issues are reported and addressed at the appropriate level promptly.
				\item To reduce transparency, as escalation limits information flow and increases confidentiality.
				\item To bypass governance, as escalation circumvents controls and standard procedures.
			\end{enumerate}
		\end{block}
		\begin{exampleblock}{Answer: B}\end{exampleblock}
		\begin{block}{Explanation}
			Escalation paths ensure timely, appropriate responses. Options A, C, and D are contrary.
		\end{block}
	\end{frame}
	
	\begin{frame}{Q629 --- AI Risk Aggregation}
		\begin{block}{Question}
			What does AI risk aggregation mean?
			\begin{enumerate}
				\item Averaging individual risks, as aggregation is a simple mean calculation and arithmetic average.
				\item Combining individual AI risks into an enterprise-wide view, accounting for interdependencies and correlations.
				\item Ignoring enterprise risk, as aggregation focuses on individual risks and ignores portfolio effects.
				\item Removing individual risks, as aggregation eliminates individual risk and reduces total exposure.
			\end{enumerate}
		\end{block}
		\begin{exampleblock}{Answer: B}\end{exampleblock}
		\begin{block}{Explanation}
			Aggregation considers interdependencies. Options A, C, and D misstate.
		\end{block}
	\end{frame}
	
	\begin{frame}{Q630 --- AI Governance Success Metrics}
		\begin{block}{Question}
			Which is a good metric for AI governance success?
			\begin{enumerate}
				\item Number of meetings, as meeting frequency correlates with governance effectiveness and organisational alignment.
				\item Reduction in AI incidents, improvements in validation timeliness, and increasing compliance with governance standards.
				\item Number of documents, as documentation volume correlates with governance effectiveness and process maturity.
				\item Cost of IT, as IT spend correlates with governance effectiveness and operational efficiency.
			\end{enumerate}
		\end{block}
		\begin{exampleblock}{Answer: B}\end{exampleblock}
		\begin{block}{Explanation}
			Governance success metrics track incidents, timeliness, and compliance. Options A, C, and D are not meaningful metrics.
		\end{block}
	\end{frame}
	
	% =================================================================
	\section{Ethics and Regulation in Depth}
	% =================================================================
	
	\begin{frame}{Q631 --- GDPR Data Subject Access Rights}
		\begin{block}{Question}
			Under GDPR, data subjects can request:
			\begin{enumerate}
				\item Only access to data, as GDPR grants a single right and is limited in scope.
				\item Access, rectification, erasure, restriction, portability, and objection.
				\item Only rectification, as GDPR grants a single right and is limited in scope.
				\item Nothing, as GDPR grants no rights to data subjects and is purely a compliance burden.
			\end{enumerate}
		\end{block}
		\begin{exampleblock}{Answer: B}\end{exampleblock}
		\begin{block}{Explanation}
			GDPR grants multiple rights. Options A, C, and D understate or misstate.
		\end{block}
	\end{frame}
	
	\begin{frame}{Q632 --- GDPR Processor vs Controller}
		\begin{block}{Question}
			What is the difference between a data controller and a data processor under GDPR?
			\begin{enumerate}
				\item They are the same, as both roles are interchangeable and functionally equivalent.
				\item Controller determines the purposes and means of processing; processor processes data on behalf of the controller.
				\item Controller only stores data, as the controller's role is limited to storage and retention.
				\item Processor only collects data, as the processor's role is limited to collection and ingestion.
			\end{enumerate}
		\end{block}
		\begin{exampleblock}{Answer: B}\end{exampleblock}
		\begin{block}{Explanation}
			Controller decides; processor acts on behalf. Options A, C, and D misstate.
		\end{block}
	\end{frame}
	
	\begin{frame}{Q633 --- GDPR Data Protection Officer}
		\begin{block}{Question}
			What is the DPO's role under GDPR?
			\begin{enumerate}
				\item Marketing, as the DPO is responsible for data-driven marketing and customer acquisition.
				\item Independent oversight of data protection compliance, advising on GDPR obligations, and liaising with the supervisory authority.
				\item Sales, as the DPO is responsible for data-driven sales and revenue growth.
				\item IT support, as the DPO is responsible for IT infrastructure and technical operations.
			\end{enumerate}
		\end{block}
		\begin{exampleblock}{Answer: B}\end{exampleblock}
		\begin{block}{Explanation}
			The DPO provides independent compliance oversight. Options A, C, and D are unrelated.
		\end{block}
	\end{frame}
	
	\begin{frame}{Q634 --- Data Protection Impact Assessment}
		\begin{block}{Question}
			What is a DPIA?
			\begin{enumerate}
				\item A marketing report, as DPIAs are used for brand positioning and competitive differentiation.
				\item A Data Protection Impact Assessment required for processing likely to result in high risk to data subjects, documenting risks and mitigations.
				\item A code review, as DPIAs are technical artefacts and code-quality assessments.
				\item A financial audit, as DPIAs are financial controls and SOX compliance mechanisms.
			\end{enumerate}
		\end{block}
		\begin{exampleblock}{Answer: B}\end{exampleblock}
		\begin{block}{Explanation}
			DPIA documents high-risk processing and mitigations. Options A, C, and D misstate.
		\end{block}
	\end{frame}
	
	\begin{frame}{Q635 --- AI Act High-Risk Systems}
		\begin{block}{Question}
			Which AI systems are classified as high-risk under the EU AI Act?
			\begin{enumerate}
				\item All chatbots, as chatbots are inherently high-risk and subject to strict obligations.
				\item Those affecting safety, fundamental rights, or critical infrastructure, such as biometrics, employment, credit scoring, and law enforcement.
				\item Only medical devices, as medical devices are the sole high-risk category and are subject to strict obligations.
				\item Only toys, as toys are the sole high-risk category and are subject to strict obligations.
			\end{enumerate}
		\end{block}
		\begin{exampleblock}{Answer: B}\end{exampleblock}
		\begin{block}{Explanation}
			High-risk systems affect safety, rights, or critical infrastructure. Options A, C, and D are narrower or incorrect.
		\end{block}
	\end{frame}
	
	\begin{frame}{Q636 --- AI Act Transparency Obligations}
		\begin{block}{Question}
			What transparency obligations does the EU AI Act impose?
			\begin{enumerate}
				\item None, as the AI Act does not address transparency and is limited to safety obligations.
				\item Notifying users when interacting with AI systems, disclosing deepfakes, and providing information about AI-generated content.
				\item Only financial disclosure, as the AI Act is limited to financial transparency and investor reporting.
				\item Only marketing disclosure, as the AI Act is limited to marketing transparency and advertising standards.
			\end{enumerate}
		\end{block}
		\begin{exampleblock}{Answer: B}\end{exampleblock}
		\begin{block}{Explanation}
			The AI Act imposes transparency obligations. Options A, C, and D are incorrect.
		\end{block}
	\end{frame}
	
	\begin{frame}{Q637 --- AI Act General Purpose AI}
		\begin{block}{Question}
			What does the EU AI Act require of general-purpose AI (GPAI) models?
			\begin{enumerate}
				\item Nothing, as GPAI models are exempt from the AI Act and are subject to no obligations.
				\item Documentation, transparency on training data, copyright compliance, and safety measures, with stricter obligations for models with systemic risk.
				\item Only marketing disclosure, as GPAI obligations are limited to marketing transparency and advertising standards.
				\item Only to register, as GPAI obligations are limited to registration and notification requirements.
			\end{enumerate}
		\end{block}
		\begin{exampleblock}{Answer: B}\end{exampleblock}
		\begin{block}{Explanation}
			GPAI models have documentation, transparency, and safety obligations. Options A, C, and D understate or misstate.
		\end{block}
	\end{frame}
	
	\begin{frame}{Q638 --- AI Act Timeline}
		\begin{block}{Question}
			How is the EU AI Act being implemented over time?
			\begin{enumerate}
				\item All at once, as the AI Act has no transitional periods and applies immediately.
				\item In phases: prohibitions first, then high-risk rules, then remaining provisions, with transitional periods for existing systems.
				\item Only in 2030, as the AI Act is delayed indefinitely and is subject to ongoing negotiation.
				\item Never, as the AI Act has been withdrawn and is no longer applicable.
			\end{enumerate}
		\end{block}
		\begin{exampleblock}{Answer: B}\end{exampleblock}
		\begin{block}{Explanation}
			The AI Act phases in obligations. Options A, C, and D are inaccurate.
		\end{block}
	\end{frame}
	
	\begin{frame}{Q639 --- CCPA Consumer Rights}
		\begin{block}{Question}
			What rights does CCPA grant California consumers?
			\begin{enumerate}
				\item No rights, as CCPA grants no consumer rights and is purely a compliance burden.
				\item Right to know, delete, opt out of sale, correct, and non-discrimination.
				\item Only access, as CCPA grants a single right and is limited in scope.
				\item Only deletion, as CCPA grants a single right and is limited in scope.
			\end{enumerate}
		\end{block}
		\begin{exampleblock}{Answer: B}\end{exampleblock}
		\begin{block}{Explanation}
			CCPA grants multiple rights. Options A, C, and D are incomplete.
		\end{block}
	\end{frame}
	
	\begin{frame}{Q640 --- PIPL Overview}
		\begin{block}{Question}
			What is the purpose of China's Personal Information Protection Law (PIPL)?
			\begin{enumerate}
				\item To ban personal data, as PIPL prohibits all personal data processing and is highly restrictive.
				\item To protect personal data, requiring consent, data minimisation, transparency, and specific rules for cross-border transfer.
				\item Only for government data, as PIPL applies solely to government data and public-sector processing.
				\item Only for healthcare data, as PIPL applies solely to healthcare data and medical records.
			\end{enumerate}
		\end{block}
		\begin{exampleblock}{Answer: B}\end{exampleblock}
		\begin{block}{Explanation}
			PIPL protects personal data with consent, minimisation, and transfer rules. Options A, C, and D are inaccurate.
		\end{block}
	\end{frame}
	
	% =================================================================
	\section{Advanced Model Validation}
	% =================================================================
	
	\begin{frame}{Q641 --- Model Validation Evidence}
		\begin{block}{Question}
			What evidence should model validation produce?
			\begin{enumerate}
				\item None, as validation is an informal process and does not require documentation.
				\item Documented assessment of conceptual soundness, data appropriateness, implementation, performance, and limitations, with findings and recommendations.
				\item Only accuracy metrics, as validation is limited to performance measurement and benchmark comparisons.
				\item Only source code, as validation is limited to code review and static analysis.
			\end{enumerate}
		\end{block}
		\begin{exampleblock}{Answer: B}\end{exampleblock}
		\begin{block}{Explanation}
			Validation produces documented, comprehensive evidence. Options A, C, and D are insufficient.
		\end{block}
	\end{frame}
	
	\begin{frame}{Q642 --- Benchmarking in Validation}
		\begin{block}{Question}
			Why does validation include benchmarking against challenger models?
			\begin{enumerate}
				\item To confuse stakeholders, as benchmarking introduces unnecessary complexity and obscures decision-making.
				\item To compare performance against alternatives, identifying whether the model is the best choice for its purpose.
				\item To reduce cost, as benchmarking reduces validation costs and increases efficiency.
				\item To speed up training, as benchmarking accelerates model development and hyperparameter search.
			\end{enumerate}
		\end{block}
		\begin{exampleblock}{Answer: B}\end{exampleblock}
		\begin{block}{Explanation}
			Benchmarking compares models for suitability. Options A, C, and D are incorrect.
		\end{block}
	\end{frame}
	
	\begin{frame}{Q643 --- Sensitivity Analysis}
		\begin{block}{Question}
			What does sensitivity analysis in model validation examine?
			\begin{enumerate}
				\item Only cost, as sensitivity analysis is limited to financial considerations and FinOps unit economics.
				\item How model outputs respond to changes in input assumptions, parameters, or data, evaluating robustness.
				\item Only speed, as sensitivity analysis is limited to latency considerations and inference-time optimisation.
				\item Only storage, as sensitivity analysis is limited to infrastructure considerations and data-lake tiering.
			\end{enumerate}
		\end{block}
		\begin{exampleblock}{Answer: B}\end{exampleblock}
		\begin{block}{Explanation}
			Sensitivity analysis tests robustness to changes. Options A, C, and D are unrelated.
		\end{block}
	\end{frame}
	
	\begin{frame}{Q644 --- Stress Testing AI Models}
		\begin{block}{Question}
			What is the purpose of stress testing AI models?
			\begin{enumerate}
				\item To make them faster, as stress testing optimizes performance and reduces latency.
				\item To evaluate their behaviour under extreme, adverse, or unlikely scenarios.
				\item To reduce cost, as stress testing reduces infrastructure costs and increases efficiency.
				\item To simplify them, as stress testing reduces model complexity and increases interpretability.
			\end{enumerate}
		\end{block}
		\begin{exampleblock}{Answer: B}\end{exampleblock}
		\begin{block}{Explanation}
			Stress tests evaluate behaviour under extremes. Options A, C, and D are unrelated.
		\end{block}
	\end{frame}
	
	\begin{frame}{Q645 --- Backtesting AI Models}
		\begin{block}{Question}
			What does backtesting involve for AI models?
			\begin{enumerate}
				\item Forward-looking simulation only, as backtesting is a forecasting technique and predictive-modelling approach.
				\item Testing the model against historical data to evaluate how it would have performed.
				\item Marketing analysis, as backtesting is a promotional activity and brand-positioning exercise.
				\item Cost analysis, as backtesting is a financial control and SOX compliance mechanism.
			\end{enumerate}
		\end{block}
		\begin{exampleblock}{Answer: B}\end{exampleblock}
		\begin{block}{Explanation}
			Backtesting evaluates historical performance. Options A, C, and D are unrelated.
		\end{block}
	\end{frame}
	
	\begin{frame}{Q646 --- Outcome Analysis}
		\begin{block}{Question}
			Outcome analysis in model validation is used to:
			\begin{enumerate}
				\item Reduce costs, as outcome analysis is a cost optimization technique and FinOps unit-economics lever.
				\item Compare model predictions against realised outcomes to detect bias, degradation, or incorrect assumptions.
				\item Speed up training, as outcome analysis accelerates model development and hyperparameter search.
				\item Increase accuracy, as outcome analysis directly improves model accuracy and benchmark performance.
			\end{enumerate}
		\end{block}
		\begin{exampleblock}{Answer: B}\end{exampleblock}
		\begin{block}{Explanation}
			Outcome analysis compares predictions to reality. Options A, C, and D are unrelated.
		\end{block}
	\end{frame}
	
	\begin{frame}{Q647 --- Use Test Evidence}
		\begin{block}{Question}
			What evidence should the use test provide?
			\begin{enumerate}
				\item None, as the use test is an informal process and does not require documentation.
				\item Qualitative evidence that the model is used appropriately, understood by users, and aligned with business decisions.
				\item Only accuracy metrics, as the use test is limited to performance measurement and benchmark comparisons.
				\item Only cost data, as the use test is limited to financial considerations and FinOps unit economics.
			\end{enumerate}
		\end{block}
		\begin{exampleblock}{Answer: B}\end{exampleblock}
		\begin{block}{Explanation}
			The use test produces qualitative evidence. Options A, C, and D are insufficient.
		\end{block}
	\end{frame}
	
	\begin{frame}{Q648 --- Model Use Restrictions}
		\begin{block}{Question}
			How should model use restrictions be enforced?
			\begin{enumerate}
				\item Through documentation only, as documentation is sufficient to enforce restrictions and ensure compliance.
				\item Through documentation, access controls, system limits, and monitoring.
				\item By ignoring them, as use restrictions are optional and impede innovation.
				\item By removing them, as use restrictions impede innovation and slow deployment.
			\end{enumerate}
		\end{block}
		\begin{exampleblock}{Answer: B}\end{exampleblock}
		\begin{block}{Explanation}
			Use restrictions require documentation, controls, limits, and monitoring. Options A, C, and D are insufficient or contrary.
		\end{block}
	\end{frame}
	
	\begin{frame}{Q649 --- Model Validation Documentation}
		\begin{block}{Question}
			What documentation should accompany model validation?
			\begin{enumerate}
				\item Nothing, as validation is an informal process and does not require documentation.
				\item A validation report with methodology, findings, limitations, and recommendations, plus supporting analyses.
				\item Only the model code, as validation is limited to code review and static analysis.
				\item Only marketing material, as validation is a promotional activity and brand-positioning exercise.
			\end{enumerate}
		\end{block}
		\begin{exampleblock}{Answer: B}\end{exampleblock}
		\begin{block}{Explanation}
			Validation documentation is comprehensive. Options A, C, and D are insufficient or inappropriate.
		\end{block}
	\end{frame}
	
	\begin{frame}{Q650 --- Responsible AI Closing Principle}
		\begin{block}{Question}
			Which best summarises the guiding principle for responsible AI governance?
			\begin{enumerate}
				\item Maximise complexity, as complexity improves governance outcomes and increases robustness.
				\item Balance performance with governance, validation, monitoring, fairness, explainability, and ethical alignment, consistent with regulatory and business expectations.
				\item Deploy without validation, as validation slows innovation and increases time-to-market.
				\item Ignore risks, as risk-taking drives innovation and increases competitive advantage.
			\end{enumerate}
		\end{block}
		\begin{exampleblock}{Answer: B}\end{exampleblock}
		\begin{block}{Explanation}
			Responsible AI balances performance and governance across all dimensions. Options A, C, and D contradict the principle.
		\end{block}
	\end{frame}
	
\end{document}